% ===========================================================================
\section{Study 2: four open-weight judges, four genres}
\label{sec:openweight}

Study~1 has three limits: 38 prompts of one genre, closed models whose
probabilities we cannot read, and vendors that may mark their text. Study~2
removes all three.

\subsection{Design}
\paragraph{Panel.} Four open-weight instruction models from four labs:
Qwen2.5-7B (Alibaba), Llama-3.1-8B (Meta), Mistral-7B-v0.3 (Mistral AI) and
Phi-4-mini (Microsoft), each run in 4-bit on one consumer GPU
(Appendix~\ref{app:openweight}). With four candidates, chance is 25\%. As in Study~1, the same
weights write each model's share of the corpus and judge it.

\paragraph{Corpus.} 120 content-scaffolded prompts in four genres: the 40
argumentative prose prompts of the original bank and 20 new ones in the same
format; 20 Python functions with fully specified behaviour; 20 explanatory
questions answered in two to four sentences; and 20 word problems with a
numeric answer. Each model answered every prompt (temperature 0.7), giving
\OWNResponses{} texts\OnPolSecondSample{}. We generated the texts ourselves
from public weights, so none carries a watermark.

\paragraph{Conditions.} The \emph{lineup} and \emph{single-text} formats of
Study~1, with the panel's names as candidates. A \emph{yes/no} format asks a
judge of every text, in a fresh context, ``Did you write this text?''; we
read the probability of \emph{yes} from the judge's next-token distribution,
so every judge answers for itself on every text and both hit and false-alarm
rates are defined. \OnPolDesign{} Every judgment is greedy.
Finally, each judge's \emph{likelihood} of every text, conditioned on the
prompt, is read directly, and each judge rates every text in both corpora
for quality.

\subsection{Results}
\begin{table*}[t]
\centering\small
\setlength\tabcolsep{3.5pt}
\begin{tabular}{l rrrrr rrr rr}
\toprule
 & \multicolumn{5}{c}{Lineup} & \multicolumn{3}{c}{Yes/no, fresh context} & \multicolumn{2}{c}{Own likelihood} \\
\cmidrule(lr){2-6}\cmidrule(lr){7-9}\cmidrule(lr){10-11}
Judge & Self & Peer & FA & Self-adv. [95\% CI] & Hit$-$FA & Hit & FA & AUC & Pick-self & AUC \\
\midrule
Qwen & 56.7\% & 15.8\% & 49.2\% & +40.8$^{*}$ [+31.1, +50.6] & +7.5 & 51.7\% & 58.6\% & 0.45 & 100.0\% & 1.00 \\
Llama & 32.5\% & 22.2\% & 36.1\% & +10.3$^{*}$ [+1.9, +18.6] & $-$3.6 & 85.8\% & 85.6\% & 0.53 & 100.0\% & 1.00 \\
Mistral & 25.0\% & 20.8\% & 20.3\% & +4.2 [$-$6.1, +14.4] & +4.7 & 99.2\% & 99.4\% & 0.41$^{*}$ & 100.0\% & 1.00 \\
Phi & 20.8\% & 22.8\% & 22.5\% & $-$1.9 [$-$10.6, +6.9] & $-$1.7 & 4.2\% & 1.4\% & 0.53 & 100.0\% & 1.00 \\
\bottomrule
\end{tabular}
\caption{Study~2, open-weight judges, all four genres pooled (120 prompts). Lineup: self-naming, peer-baseline and false-alarm (FA) rates, the self-advantage with a prompt-bootstrap interval, and the hit minus false-alarm rate, which is positive only if a judge names itself more on its own text than on others'. Yes/no: ``Did you write this text?'' in a fresh context; hit and false-alarm rates at $P(\text{yes})>0.5$ and the AUC of $P(\text{yes})$ for own versus others' text (0.5 is chance). Own likelihood: the share of prompts on which the text the judge finds most likely, relative to the other judges, is its own (chance 25\%), and the AUC of that relative likelihood. $^{*}$Holm-corrected within-prompt permutation $p<0.05$ across the four judges.}
\label{tab:openweight}
\end{table*}


\paragraph{No open-weight judge recognises itself in the lineup.} Two
judges clear the first test of Section~\ref{sec:measure} and fail the
second (Table~\ref{tab:openweight}). Qwen names itself on
\OWQwenLineupSelf{} of its own texts while its peers name Qwen on only
\OWQwenLineupPeer{}, a self-advantage of \OWQwenLineupAdv{} points
\OWQwenLineupAdvCI{}; but it also names itself on \OWQwenLineupFA{} of the
texts it did not write, so its hits exceed its false alarms by only
\OWQwenLineupDisc{} points (Holm $p=$~\OWQwenLineupDiscHolm{}). Llama's
self-advantage of \OWLlamaLineupAdv{} \OWLlamaLineupAdvCI{} comes with
\emph{more} false alarms than hits (\OWLlamaLineupFA{} against
\OWLlamaLineupSelf{}). Both judges beat their peers on their own text only
because they put their own name on much of what they read. Mistral and Phi show neither
an advantage nor discrimination. The judges attribute the lineup at
\OWLineupAcc{} overall, against 25\% chance, although the texts carry the
signal: the classifier of Section~\ref{sec:frontier}, trained and tested on
disjoint prompts, attributes them at \OWClfPooled{} \OWClfPooledCI{}, from
\OWClfMath{} on math to \OWClfProse{} on prose.

\paragraph{One text at a time, each judge answers by habit.} The frontier
judges, shown one text at a time, piled their guesses onto GPT and Claude.
The open-weight judges pile them onto themselves, or onto no one: Qwen names
itself on \OWQwenSingleNamesSelf{} of all texts and Llama on
\OWLlamaSingleNamesSelf{}, whoever wrote them, while Mistral names itself on
only \OWMistralSingleNamesSelf{}. The self-advantage then measures the habit.
It is \OWQwenSingleAdv{} points for Qwen, \OWLlamaSingleAdv{} for Llama,
\OWPhiSingleAdv{} for Phi and \OWMistralSingleAdv{} for Mistral, all four
significant after Holm correction, and none of the four judges
discriminates (hit minus false-alarm rate \OWLlamaSingleDisc{} to
\OWPhiSingleDisc{} points, Holm $p\geq$~\OWPhiSingleDiscHolm{}). Read as
self-naming rates, this format would credit Qwen and Llama with near-perfect
self-recognition.

\paragraph{Asked directly, judges answer by disposition.} The yes/no format
removes the candidate list, so a judge cannot pile its answers onto one
popular name; it can only say yes or no to each text. What it says depends
almost entirely on the judge (Table~\ref{tab:openweight}). Mistral says yes
to \OWMistralBinaryHit{} of its own texts and to \OWMistralBinaryFA{} of
everyone else's; Llama to \OWLlamaBinaryHit{} and \OWLlamaBinaryFA{}; Qwen
to \OWQwenBinaryHit{} and \OWQwenBinaryFA{}; Phi to \OWPhiBinaryHit{} and
\OWPhiBinaryFA{}. Read as self-naming rates, these would rank the four judges
from near-certain self-recognition to none; read against the false alarms,
none of them discriminates. The AUC of $P(\text{yes})$, which is free of that
disposition, is \OWQwenBinaryAUC{} for Qwen, \OWLlamaBinaryAUC{} for Llama,
\OWPhiBinaryAUC{} for Phi and \OWMistralBinaryAUC{} \OWMistralBinaryAUCCI{}
for Mistral, which is reliably \emph{below} chance (within-prompt
permutation, Holm $p=$~\OWMistralBinaryAUCHolm{}): Mistral is slightly more
willing to claim other models' text than its own. Split by genre, the
highest AUC of any judge is \OWBinaryAUCGenreMax{} \OWBinaryAUCGenreMaxCI{}
(\OWBinaryAUCGenreMaxWho{}).

\subsection{On-policy judging}
\label{sec:onpolicy}
\citet{khullar2026selfattribution} find that a monitor's judgment shifts
most when it evaluates an action it produced earlier in the same
conversation. Our on-policy conditions ask whether that position also lets a
judge recognise its writing. When the candidate is the judge's own previous
turn, word for word, every judge's $P(\text{yes})$ separates it from peers'
answers: the AUC is \OWMistralOnPolAUC{} for Mistral, \OWQwenOnPolAUC{} for
Qwen, \OWPhiOnPolAUC{} for Phi and \OWLlamaOnPolAUC{} for Llama (all Holm
$p<0.01$). That is a low bar, since the candidate is a copy of text already in
the context window, and even so it is not met everywhere: on prose Llama's
AUC is \OWLlamaOnPolProseAUC{}, below chance, and its thresholded answers keep
the habits of the fresh-context format (Mistral says yes to
\OWMistralOnPolFA{} of peers' answers). \OnPolResampled{}
